\section{Case study: a repair centre for medical equipment}\label{sec:case}

\bandeau

\subsection{Site and data}

The repair centre receives field-replaceable units of medical imaging systems. \apres{one paragraph on the site:
families of equipment, yearly volumes of repaired and scrapped units, the bulk contract and what it costs, the
people who use the system, and since when.} The system reads the centre's operational extractions as they are
exported: repair and scrap records, parts consumption, weighing of the waste bins, a file of early failures of
installed parts, and the bill of materials. Parameters that no extraction contains (operator times, values of
reused parts, recycler prices) are entered in a workbook as ranges with their provenance; estimates are labelled as
such, and a parameter left empty stays unknown. Unit types and components are not named in this paper.
\noteechelle

The line model covers \CasTypes{} unit types, of which \CasIncomplets{} could not be evaluated for lack of
parameters and are reported as such rather than counted as zero. The pipeline ended with \CasAnomalies{} failed
computations. \apres{the verification of the estimated input flow against the weighed mass: number of weeks
compared, share within the stated margin, and what was corrected.}

\subsection{The decision}

Table~\ref{tab:case-decision} gives the indicators on which the decision to install the line rests, at the centre of
the parameter ranges, and the distribution of the annual result when the parameters are drawn from their ranges.
The optimal plan splits \CasPartages{} unit type between two plans, in line with Proposition~\ref{prop:split}.
\apres{which unit types are opened and which are left, which components are reused and which are sorted, and
whether the plan matches what the technicians expected; who validated the valuation of reused parts.}

\begin{table}[H]
\centering\small
\caption{Case study: decision indicators. \ifnum\DonneesDemo=1\relax\textbf{Demonstration data, not
results.}\fi}\label{tab:case-decision}
\begin{tabular}{lr}
\toprule
Indicator & Value \\
\midrule
Variable margin of the line (per year) & \CasMarge \\
Fixed costs of the line, annualised & \CasFixes \\
Annual result against the current contract, $V-K$ & \CasResultat \\
Current cost of bulk treatment (per year) & \CasCoutActuel \\
Operator hours used / available (per year) & \CasHeures{} / \CasCapacite \\
Value of one more operator hour, $\mu^\ast$ & \CasMu \\
Units processed (per year) & \CasUnites \\
Break-even volume (units per year) & \CasSeuil \\
Payback (years) & \CasRetour \\
Mass diverted from the bulk stream / mass scrapped (kg per year) & \CasDetournee{} / \CasMasse{} (\CasPartDetournee) \\
\midrule
Probability that the line is profitable (parameter ranges) & \CasProba \\
Annual result: 5th percentile / median / 95th percentile & \CasPcinq{} / \CasMediane{} / \CasPnonantecinq \\
\bottomrule
\end{tabular}
\end{table}

Table~\ref{tab:case-hours} varies the operator hours. It displays Proposition~\ref{prop:dual} on the site's data:
the margin is concave in the hours, and the value of an additional hour falls by steps until the hours cease to
bind. \apres{the staffing decision this table supported.}

\begin{table}[H]
\centering\small
\caption{Case study: sensitivity to operator hours. \ifnum\DonneesDemo=1\relax\textbf{Demonstration
data.}\fi}\label{tab:case-hours}
\begin{tabular}{rrrrr}
\toprule
Hours per year & Variable margin & Annual result & Occupancy & Value of one more hour \\
\midrule
\corps{t_cas_heures}
\bottomrule
\end{tabular}
\end{table}

\subsection{Is a dynamic policy worth building?}

Replayed week by week on the \CasPeriodes{} weeks of history, with the fixed cost paid every week and no stock, the
line earns \CasMyope{} per year, against \CasFluide{} for the model on the average flow. The variability gap of
Proposition~\ref{prop:fluid} is therefore \CasGamma{} per year, \CasGammaRel{} of the fluid value;
\CasPertes{} weeks end at a loss. \apres{the conclusion drawn: whether the gap justifies a buffer of scrapped
units, and what was decided.}

\subsection{Yields}

The posterior on yields was scored prequentially on \CasLots{} teardown batches (\CasCopies{} copies tested).
Table~\ref{tab:case-yields} compares it with the prior alone and with the raw observed frequency. \apres{which
priors the audit corrected most, in which direction, and whether the three biases named in
Section~\ref{sec:learning} showed.}

\begin{table}[H]
\centering\small
\caption{Case study: prequential scores of the yield estimate (lower is better, except coverage, whose target is
90\%). \ifnum\DonneesDemo=1\relax\textbf{Demonstration data.}\fi}\label{tab:case-yields}
\begin{tabular}{lrrrr}
\toprule
Estimate & Brier score & Log loss & Absolute error & Coverage of the 90\% interval \\
\midrule
\corps{t_cas_rendements}
\bottomrule
\end{tabular}
\end{table}

\subsection{Arrivals}

Table~\ref{tab:case-forecast} compares the forecasting methods of Section~\ref{sec:forecast} by rolling origin on
the weekly series of scrapped units. \apres{the reading of the table; whether the structural model, which uses
the inflow, beats the univariate methods; and the result of any zero-shot foundation model run on the same
series.}

\begin{table}[H]
\centering\small
\caption{Case study: forecasts of weekly scrapped units, rolling origin, all methods scored on the same
forecasts. Root mean squared error relative to the historical mean; pinball loss averaged over quantiles; coverage
of the 80\% interval. \ifnum\DonneesDemo=1\relax\textbf{Demonstration data.}\fi}\label{tab:case-forecast}
\begin{tabular}{lrrrr}
\toprule
Method & Forecasts & Relative RMSE & Pinball loss & Coverage (80\%) \\
\midrule
\corps{t_cas_previsions}
\bottomrule
\end{tabular}
\end{table}

\subsection{Composition of components}

The bill of materials of the modelled units lists \CasCompoPieces{} components. After the run, \CasCompoValidees{}
have a class validated by a human, \CasCompoAcceptees{} a class accepted without review and \CasCompoFile{} wait
in the review queue. The chain made \CasCompoAppels{} calls to the language model in this run, at a cost of
\CasCompoCout; \CasCompoSansAppel{} components were handled without any paid acquisition.
\apres{the model used and where it ran; the measured reliability of each source type after the audit; the number of
reviews done and the time they took; the result of the instruction search on the validated parts (held-out
accuracy before and after, with its test); and one or two examples, anonymised, of components whose class changed
a decision.}

\subsection{What the agents produced}

\apres{the decision note written by the agents on the site's data: how many tool calls the pilot made, how many
numbers the grounding check flagged, the reviewer's reservations, how the note compares with the one a project
engineer wrote, and what the manager did with it.}
